% 51-40(51쪽) — 포물선 y=2x^2-6x+7 과 직선 y=2x-6, 직선 위의 두 점 A(4,2)·B(6,6), 포물선 위의 점 P 로 만든 삼각형 ABP. 스캔 화질이 낮아 TikZ 로 다시 그림.
% 원문에 있는 것만: 눈금 숫자 7 6 2 (y) · 4 6 (x) · 라벨 O, x, y, P, A, B, y=2x^2-6x+7, y=2x-6 · 좌표 안내 파선 · 회색 삼각형.
% P 는 「임의의 점」이라 원문이 찍은 자리(x≒2.35)에 그대로 두어 넓이가 최소가 되는 자리를 그림에서 읽을 수 없게 한다. 원문에 없는 눈금은 넣지 않는다.
% 점 표시: A·B 만 채운 점(원문) · P 는 곡선 위 꼭짓점만(원문에 점 표시 없음).
\begin{tikzpicture}[x=0.40cm, y=0.36cm, line width=0.5pt]
  \draw[->, >=stealth, line width=0.7pt] (-1.7,0) -- (8.2,0) node[below=1pt] {$x$};
  \draw[->, >=stealth, line width=0.7pt] (0,-2.2) -- (0,10.2) node[left=1pt] {$y$};
  \node[below left=1pt and -1pt, font=\small] at (0,0) {$\pt{O}$};
  % 좌표 안내 파선(원문)
  \draw[dashed, line width=0.35pt] (0,6) -- (6,6);
  \draw[dashed, line width=0.35pt] (0,2) -- (4,2);
  \draw[dashed, line width=0.35pt] (4,0) -- (4,2);
  \draw[dashed, line width=0.35pt] (6,0) -- (6,6);
  % 삼각형 ABP(회색 칠 · 원문)
  \fill[black!15] (2.35,3.945) -- (4,2) -- (6,6) -- cycle;
  \draw[line width=0.5pt] (4,2) -- (2.35,3.945) -- (6,6);
  % 포물선과 직선
  \draw[line width=0.6pt, domain=-0.33:3.33, samples=120, smooth] plot (\x, {2*\x*\x-6*\x+7});
  \draw[line width=0.6pt] (1.75,-2.5) -- (6.75,7.5);
  % 점
  \fill (4,2) circle (2.2pt);
  \fill (6,6) circle (2.2pt);
  % 라벨(원문 자리 그대로: 7 은 축 오른쪽, 6·2 는 축 왼쪽)
  \node[left=1pt, font=\small] at (2.35,3.945) {$\pt{P}$};
  \node[right=3pt, font=\small] at (4,2) {$\pt{A}$};
  \node[right=3pt, font=\small] at (6,6) {$\pt{B}$};
  \node[right=2pt, font=\small] at (0,7) {$7$};
  \node[left=2pt, font=\small] at (0,6) {$6$};
  \node[left=2pt, font=\small] at (0,2) {$2$};
  \node[below=2pt, font=\small] at (4,0) {$4$};
  \node[below=2pt, font=\small] at (6,0) {$6$};
  \node[anchor=west, font=\small] at (0.55,9.6) {$y=2x^2-6x+7$};
  \node[anchor=west, font=\small] at (4.3,8.3) {$y=2x-6$};
\end{tikzpicture}
